% XPtoIRL installation and user manual (compile with XeLaTeX)
\documentclass[a4paper,11pt]{article}
\usepackage[top=2cm, bottom=2cm, left=2cm, right=2cm]{geometry}
\usepackage{fontspec}
\setmainfont{Arial}
\setmonofont{Consolas}[Scale=0.9]
\usepackage{polyglossia}
\setdefaultlanguage{english}
\usepackage{setspace}
\setstretch{1.2}
\setlength{\parskip}{6pt}
\setlength{\parindent}{0pt}
\usepackage{titlesec}
\titleformat{\section}{\fontsize{13}{15}\selectfont\bfseries}{\thesection}{0.6em}{}
\titleformat{\subsection}{\fontsize{11.5}{14}\selectfont\bfseries}{\thesubsection}{0.6em}{}
\usepackage{enumitem}
\setlist{itemsep=2pt, topsep=2pt}
\usepackage{fancyvrb}
\usepackage{xcolor}
\usepackage{tabularx}
\usepackage{array}
\newcolumntype{L}[1]{>{\raggedright\arraybackslash}p{#1}}
\usepackage[hidelinks]{hyperref}
\RecustomVerbatimEnvironment{Verbatim}{Verbatim}{frame=single, rulecolor=\color{gray!60}, xleftmargin=4pt}

\begin{document}

{\fontsize{16}{19}\selectfont\bfseries XPtoIRL: Installation and User Manual}\\[2pt]
Release 1.0.0 \textbar\ Project: Software Engineering (DLMCSPSE01) \textbar\ Leela Sai Abhinav Gamidi (UPS10617430)

\section{Overview}
XPtoIRL is a web application that recommends video games for the real-life skills a
user wants to develop. The application can be used in three ways: online in the cloud
(no installation), locally with Docker (recommended), or locally with Python only.

\section{Online use}
Open the link given in the submission file \texttt{XPtoIRL\_links.txt} in a current
browser (Chrome, Firefox or Edge). The file also contains the login data of a test
administrator. On the free cloud plan the first request after a longer pause can take
up to one minute, because the service is started again.

\section{Installation with Docker (recommended)}
\textbf{Requirements:} Docker Desktop with Docker Compose, Git.

\begin{Verbatim}
git clone <repository URL>
cd xptoirl
docker compose up --build
\end{Verbatim}

Open \url{http://localhost:8000}. On the first start the container creates the
PostgreSQL tables and imports the data set (17 skills, 66 games, 175 mappings,
26 sources). If port 8000 is already used, choose another host port:

\begin{Verbatim}
# macOS / Linux
WEB_PORT=8010 docker compose up --build
# Windows PowerShell
$env:WEB_PORT=8010; docker compose up --build
\end{Verbatim}

To create an administrator, set \texttt{ADMIN\_EMAIL} and \texttt{ADMIN\_PASSWORD} in a
\texttt{.env} file (copy \texttt{.env.example}) before the start, or run:

\begin{Verbatim}
docker compose exec web flask create-admin admin@example.org
\end{Verbatim}

Stop the application with \texttt{docker compose down}; add \texttt{-v} to delete the database as well.

\section{Installation without Docker}
\textbf{Requirements:} Python 3.11 or newer. Without \texttt{DATABASE\_URL} the application uses a local SQLite file.

\begin{Verbatim}
python -m venv .venv
.venv\Scripts\activate          # macOS/Linux: source .venv/bin/activate
pip install -r requirements-dev.txt
set FLASK_APP=wsgi.py           # macOS/Linux: export FLASK_APP=wsgi.py
flask db upgrade
flask seed
flask create-admin admin@example.org
flask run
\end{Verbatim}

Open \url{http://127.0.0.1:5000}.

\section{Cloud deployment (Render)}
The repository contains the Render Blueprint \texttt{render.yaml}. In the Render
dashboard choose \textbf{New $>$ Blueprint}, select the GitHub repository, enter
\texttt{ADMIN\_EMAIL} and \texttt{ADMIN\_PASSWORD} and confirm. Render creates the
PostgreSQL database and the web service, generates \texttt{SECRET\_KEY}, sets
\texttt{DATABASE\_URL} and checks the service with \texttt{/api/health}. Each push to
the main branch deploys the new version.

\section{Configuration}
\begin{tabularx}{\textwidth}{|l|X|l|}
\hline
\textbf{Variable} & \textbf{Meaning} & \textbf{Default} \\ \hline
\texttt{APP\_ENV} & \texttt{development}, \texttt{production} or \texttt{testing} & \texttt{development} \\ \hline
\texttt{SECRET\_KEY} & Key for sessions and CSRF tokens, required in production & dev value \\ \hline
\texttt{DATABASE\_URL} & SQLAlchemy URL of the database & SQLite file \\ \hline
\texttt{ADMIN\_EMAIL}, \texttt{ADMIN\_PASSWORD} & Administrator created on container start & empty \\ \hline
\texttt{WEB\_PORT} & Host port used by \texttt{docker compose} & 8000 \\ \hline
\texttt{PORT} & Port of Gunicorn inside the container (set by the cloud) & 8000 \\ \hline
\end{tabularx}

\section{Using the application}
\begin{enumerate}[leftmargin=*]
\item \textbf{Find games.} On the start page choose one to three skills (grouped by
  cognitive, social, motor and emotional) and press \emph{Find games}.
\item \textbf{Filter.} Narrow the results by platform, genre, PEGI age, price model,
  play mode and session length. Each result shows the score, the reason, the
  evidence level (A strong, B medium, C weak) and the sources.
\item \textbf{Game page.} Open a game to see all trained skills with reasons and
  links to the studies. The search field finds games by title.
\item \textbf{Learning plan.} After registering, save games to the learning plan, rate
  how well a game trained a skill (1 to 5) and print the plan. Under
  \emph{Account} the account can be deleted together with all data.
\item \textbf{Administration.} Administrators open \emph{Admin} to create, edit and
  delete skills, games, sources and mappings, or to import the four CSV files.
  Every mapping needs at least one source.
\end{enumerate}

\section{Tests and quality checks}
\begin{Verbatim}
pytest --cov=app                                   # 86 tests
ruff check .                                       # linter
python scripts/load_test.py --base-url http://localhost:8000 --users 20
\end{Verbatim}

GitHub Actions runs the linter, the tests with a coverage gate (service layer
$\geq$ 70\,\%) and a Docker build for every push.

\section{Troubleshooting}
\begin{tabularx}{\textwidth}{|L{5.2cm}|X|}
\hline
\textbf{Problem} & \textbf{Solution} \\ \hline
\emph{port is already allocated} & Start with another \texttt{WEB\_PORT} (Section 3). \\ \hline
Docker cannot connect to the daemon & Start Docker Desktop and wait until it shows ``running''. \\ \hline
\emph{SECRET\_KEY must be set in production} & Set \texttt{SECRET\_KEY} or use \texttt{APP\_ENV=development} locally. \\ \hline
Empty start page & Run \texttt{flask seed} (or restart the container) to import the data set. \\ \hline
\end{tabularx}

\end{document}
